% TOP_PROBLEM: {"id": 30006884, "title": "Benjamini-Schramm conjecture p_c<p_u for non-amenable Cayley graphs", "category_id": 19, "category": {"id": 19, "name": "probability", "display_name": "Probability", "description": "Problems involving probability theory, stochastic processes, and random structures.", "slug": "probability", "order_index": 19, "created_at": "2026-07-31T15:26:25.671Z"}, "status": "open"}
% FIRST_POSTED: 2026-09-27
% !TeX program = lualatex
% ATTEMPT_STATUS: unresolved
% Model: GPT-6 Astra Ultra; continuation dated 27 September 2026.
\documentclass[11pt]{article}
\usepackage[margin=25mm]{geometry}
\usepackage{fontspec}
\setmainfont{Latin Modern Roman}
\usepackage{amsmath,amssymb,mathtools}
\usepackage{unicode-math}
\setmathfont{Latin Modern Math}
\usepackage{xurl,hyperref}
\hypersetup{colorlinks=true,urlcolor=blue}
\setcounter{secnumdepth}{0}
\setlength{\parindent}{0pt}
\setlength{\parskip}{5pt}
\setlength{\emergencystretch}{3em}
\begin{document}
\section*{\texorpdfstring{Benjamini-Schramm conjecture p\_c<p\_u for non-amenable Cayley graphs}{Benjamini-Schramm conjecture p\_c<p\_u for non-amenable Cayley graphs}}\label{30006884}
\subsection{Definitions and mathematical statement}
For a locally finite Cayley graph $G$, Bernoulli bond percolation independently retains each edge with probability $p$. Let $N_\infty$ be the number of infinite open clusters and let $o$ be a fixed vertex. Define
\[
 p_c=\inf\{p:{\mathbb{P}}_p(o\leftrightarrow\infty)>0\},\qquad
 p_u=\inf\{p:{\mathbb{P}}_p(N_\infty=1)=1\}.
\]
A finitely generated group is amenable if it has a translation-invariant mean, equivalently a F\o lner sequence of finite sets; nonamenability is failure of this property. The selected assertion is $p_c(G)<p_u(G)$ exactly when the infinite underlying group is nonamenable, for every finite generating set defining a Cayley graph. The thresholds are infima, so the statement does not prescribe cluster behavior at either endpoint. Site percolation and arbitrary nontransitive graphs are not silently substituted.
\par\smallskip\textit{Formulation and concept references:} \href{https://ar5iv.labs.arxiv.org/html/1712.04651}{[S1]}.

\subsection{Short English statement}
Does a Cayley graph have a nonempty percolation phase with multiple infinite clusters exactly when its group is nonamenable?
\subsection{Research attempt}
\paragraph{The fixed generating set matters.}
The inspected formulation in [S1, Conjecture 2] is about every
finite generating set of the given infinite group. Its discussion
of special generating systems does not establish that quantifier.
Likewise, a planar site-percolation result or a nontransitive graph
example is not silently transferred to Bernoulli bond percolation
on an arbitrary Cayley graph.

We prove the complete regular-tree subcase, then derive an explicit
isoperimetric/spectral sufficient condition for a nonuniqueness
interval on a regular graph. The latter condition is stronger than
nonamenability. Identifying that extra quantitative requirement
pinpoints the limitation of this approach.

Assume the Cayley graph is undirected, connected and locally finite,
of degree $d$. Edges are independently open with probability $p$.
An infinite cluster contains an infinite simple ray from each of
its vertices: apply the elementary infinite-branching-path lemma
to the locally finite tree of finite simple paths in the cluster.
This observation connects infinite connectivity with the finite-path
estimates used below.

\paragraph{A universal elementary lower bound on criticality.}
The number of simple paths of length $n\geq1$ starting at $o$ is
at most $d(d-1)^{n-1}$. Each has $n$ distinct edges and is open
with probability $p^n$. The union bound gives
\[
 \mathbb P_p(o\hbox{ has an open simple path of length }n)
      \leq d p\,[p(d-1)]^{n-1}.
\]
If $p(d-1)<1$, the right side tends to zero, while an infinite
cluster would supply every such length. Therefore
\[
 p_c\geq\frac1{d-1}.
 \tag{360.1}
\]
This does not decide the endpoint, and the target only concerns
the threshold infimum. Replacing the simple-path count by an exact
tree count in a graph with cycles would not be legitimate.

\paragraph{Exact criticality on the regular tree.}
Let $\mathbb T_d$ be the $d$-regular tree with $d\geq3$. After
one oriented edge, each vertex has $d-1$ forward neighbors. Its
forward open descendants form a Galton--Watson process with
offspring distribution $\operatorname{Bin}(d-1,p)$. If $q$ is its
extinction probability, independence gives
\[
 q=(1-p+pq)^{d-1},
 \tag{360.2}
\]
where $q$ is the smallest fixed point in $[0,1]$. This last
qualification follows by iterating the offspring generating function
from zero: the $n$-th iterate is extinction by generation $n$.

For $0<p<1$, the generating function is strictly convex. If its
derivative at one, $(d-1)p$, is at most one, its graph lies strictly
above the diagonal before one, except for the irrelevant degenerate
cases; hence $q=1$. If the derivative is greater than one, the
function minus the diagonal is negative just to the left of one
and positive at zero, so it has a fixed point below one. This proves
survival exactly when $p>1/(d-1)$. At the root, the survival probability is
\[
 \theta(p)=1-(1-p+pq)^d,
 \tag{360.3}
\]
and consequently $p_c(\mathbb T_d)=1/(d-1)$.

\paragraph{Infinitely many infinite clusters for every subunit supercritical p.}
Fix $p\in(1/(d-1),1)$. Choose an infinite ray in the tree and,
at each successive vertex, choose a branch pointing off the ray.
The corresponding forward rooted subtrees are pairwise disjoint;
for $d\geq3$ there is at least one such branch at every vertex.
For each chosen subtree require its attaching edge to be closed
and its root's forward cluster to be infinite.

These events depend on pairwise disjoint edge sets. They are
independent and each has the same positive probability
$(1-p)(1-q)$. The probability that only finitely many occur is
zero: for any tail, the probability that none occurs in its first
$m$ branches is $(1-(1-p)(1-q))^m\to0$, and take a countable union
over possible last occurrences. When an event occurs, its infinite
cluster is separated from the rest of the tree by the closed
attaching edge. Distinct successful branches therefore supply
distinct infinite clusters. We have proved
\[
 N_\infty=\infty\quad\hbox{almost surely for }
          1/(d-1)<p<1.
 \tag{360.4}
\]
For $p\leq1/(d-1)$ there are no infinite clusters, by the branching
argument and a countable union over vertices. At $p=1$ the whole
tree is connected. Hence $p_u(\mathbb T_d)=1$ and
$p_c<p_u$ exactly as requested in this subcase.

The tree is nonamenable directly: for a finite nonempty vertex set
$F$, its induced forest has at most $|F|-c$ edges, where $c\geq1$
is its number of components. The edge boundary therefore has size
\[
 |\partial_EF|=d|F|-2|E(F)|\geq(d-2)|F|+2c.
\]
In particular its isoperimetric constant is positive. Even-degree
regular trees are the standard Cayley graphs of free groups, so
this yields actual instances of the catalogue statement, not just
an unrelated nontransitive branching model.

\paragraph{An isoperimetric upper bound on p_c, with exploration proof.}
For a general graph define the edge constant
\[
 h=\inf_{\varnothing\ne F\text{ finite}}
                         \frac{|\partial_EF|}{|F|}.
\]
We prove
\[
 p_c\leq\frac1{1+h}\quad\hbox{if }h>0.
 \tag{360.5}
\]
Explore the cluster of $o$ by testing an unexamined edge with one
endpoint in the discovered cluster and one outside. On an open
test add the new endpoint; on a closed test add none. Discard
unexamined edges whose endpoints have both become discovered,
since they need not be tested. After $t$ tests, let $O_t,C_t$ be
the open and closed counts. The discovered vertex set has
$1+O_t$ vertices, and the number of its still untested outgoing
edges is at least
\[
 h(1+O_t)-C_t=h+(h+1)O_t-t.
 \tag{360.6}
\]
Every already tested edge in this boundary was closed, so subtracting
$C_t$ is a valid, possibly wasteful bound. The adaptive order of
testing does not change the fact that each fresh edge is an
independent Bernoulli-$p$ variable.

Imagine an infinite independent sequence of test results, and put
$S_t=h+(h+1)O_t-t$. Its increment is $h$ on an open result and
$-1$ on a closed result. If $p>1/(h+1)$, its mean increment is
positive. This random walk has positive probability to remain
strictly positive forever. To see this without assuming a ruin
formula, the strong law makes an independent future walk tend to
$+\infty$, so its infimum is finite almost surely. For a sufficiently
large integer $K$, its infimum exceeds $-Kh$ with positive
probability. Force the first $K+1$ results to be open, an event
of positive probability independent of the future. Their reserve
then keeps $S_t$ positive at all times on that future event.

Whenever $S_t>0$, (360.6) guarantees a further edge to test, so
the actual exploration cannot stop. Infinitely many tests with the
positive open density then discover infinitely many vertices. This
proves positive survival probability and hence (360.5).

\paragraph{A spectral obstruction to uniqueness.}
Let $P$ be the simple-random-walk operator on $\ell^2(V)$,
$Pf(x)=d^{-1}\sum_{y\sim x}f(y)$, and let $\rho=\|P\|_{2\to2}$.
It is self-adjoint for an undirected regular graph. Write
$\tau_p(x,y)=\mathbb P_p(x\leftrightarrow y)$. A connection
contains a simple path, so the path union bound gives
\[
 \tau_p(x,y)\leq\sum_{n\geq0}(pd)^nP^n(x,y).
 \tag{360.7}
\]
Indeed simple paths of length $n$ are bounded in number by all
walks of length $n$, namely $d^nP^n(x,y)$. We do not assign
probability $p^n$ to a repeated-edge walk; it is only an upper
bound on the number of simple paths.

If $pd\rho<1$, the operator series on the right converges on
$\ell^2$ to $(I-pdP)^{-1}$. Its row at $o$ is an $\ell^2$
sequence, since the operator is bounded and self-adjoint. By
nonnegativity and (360.7), $\tau_p(o,\cdot)$ is also in $\ell^2$.
In particular it cannot be bounded below by a positive constant
on infinitely many vertices.

Suppose instead that there is almost surely a unique infinite
cluster. Its existence implies $\theta(p)>0$: if every vertex
had zero chance of belonging to it, a countable union would give
zero chance of any infinite cluster. Transitivity makes the chance
the same at every vertex. The increasing events that $o$ and $x$
belong to infinite clusters are positively correlated under product
measure, by the Bernoulli correlation inequality. If both occur,
uniqueness connects them. Thus
\[
 \tau_p(o,x)\geq\mathbb P_p(o\leftrightarrow\infty,
                              x\leftrightarrow\infty)
                  \geq\theta(p)^2>0
 \tag{360.8}
\]
for every $x$, a contradiction. The correlation inequality follows
first for increasing finite-coordinate events by induction on the
independent bits, then for these events by limits. We conclude
\[
 p_u\geq\frac1{d\rho}.
 \tag{360.9}
\]
For the connected infinite regular graphs under discussion $d\rho$
is greater than one; the displayed bound is used only in its
meaningful probability range.

\paragraph{A rigorous quantitative subcase of the conjecture.}
Combining (360.5) and (360.9) proves
\[
 1+h>d\rho\quad\Longrightarrow\quad
 p_c\leq\frac1{1+h}<\frac1{d\rho}\leq p_u.
 \tag{360.10}
\]
This gives a complete sufficient condition for the selected strict
inequality. Its proof establishes a positive interval, rather than
inferring strictness from two non-strict bounds with the same
endpoint.

Nonamenability supplies positive isoperimetry, but it does not give
the stronger numerical comparison $1+h>d\rho$ by any argument
above. The two estimates can leave overlapping ranges. Nor may
one change the generating set to improve $d,h,\rho$ and then
declare the original graph settled. Those parameters and both
thresholds depend on the chosen graph, whereas the requested
quantifier is over every generating set.

The tree proof illustrates a further limitation of (360.10): the
disjoint-branch argument supplies nonuniqueness directly and can
work when a crude combination of numerical bounds does not.
For a graph with cycles, closing one attaching edge need not
isolate an infinite branch, so the same construction cannot be
copied without checking all alternative connections.

\paragraph{Diagnostics and outcome.}
The finite computation iterates the branching generating function
from zero for selected $d,p$, checks the explicit fixed points
where available, and enumerates small finite subtrees to verify
the edge-boundary identity. It records the formula for criterion
(360.10), without testing parameter triples. Finite-depth
survival probabilities decrease toward the infinite survival
probability; positive values at finite depth alone do not prove
supercriticality, so the exact convexity argument remains essential.

The complete regular-tree result and the quantitative criterion are
proved. The first missing step for arbitrary nonamenable Cayley
graphs is a nonuniqueness mechanism beyond (360.10), or a stronger
bound that forces a strict interval for every generating set.
The reverse amenable direction requires the established uniqueness
theory rather than a tree computation; it is not reproved here.
The full equivalence remains unresolved by this attempt.
\subsection{Sources}
\begin{itemize}
\item[S1]H. Duminil-Copin, "Sixty years of percolation", Proc. International Congress of Mathematicians 2018, vol. IV, pp. 639-672; arXiv:1712.04651.
\url{https://ar5iv.labs.arxiv.org/html/1712.04651}.
\textit{Formulation source.}
\end{itemize}
\end{document}
